% Problem dossier: VI.02.P2
\subsection*{Problem VI.02.P2 --- Principal opens form a basis}
\label{prob:vi02-p2}

\paragraph{Problem.}
Let \(X\subseteq\mathbb A_k^n\) be an affine algebraic set.  Prove that the collection
\[
\set{D_X(f):f\in k[x_1,\dots,x_n]}
\]
is a basis for the Zariski topology on \(X\).  Also prove
\[
D_X(f)\cap D_X(g)=D_X(fg).
\]

\ifdefined\FullProblemDossiers
\paragraph{Why this problem matters.}
A topology is usable when one has manageable local neighborhoods.  Principal opens are the
canonical neighborhoods of affine algebraic geometry and later become the basic local pieces
of affine schemes.

\paragraph{Useful ideas before starting.}
Write the complement of an open set as a relative zero locus \(V_X(S)\).  A point outside
that zero locus must fail at least one equation in \(S\).

\paragraph{Guided hints.}
\textbf{Hint 1.} If \(x\in U\) and \(X\setminus U=V_X(S)\), choose \(f\in S\) with
\(f(x)\neq0\).

\textbf{Hint 2.} Show \(D_X(f)\subseteq U\).

\textbf{Hint 3.} For intersections, use that a product in a field is nonzero exactly when
both factors are nonzero.

\paragraph{Solution.}
Let \(U\subseteq X\) be open and \(x\in U\).  Since \(X\setminus U\) is closed, write
\[
X\setminus U=V_X(S)
\]
for some family \(S\).  Since \(x\notin V_X(S)\), there is \(f\in S\) with
\(f(x)\neq0\).  Thus \(x\in D_X(f)\).  If \(y\in D_X(f)\), then \(f(y)\neq0\), so
\(y\notin V_X(S)\), hence \(y\in U\).  Therefore
\[
x\in D_X(f)\subseteq U.
\]
This is precisely the basis criterion.

For the intersection formula,
\[
x\in D_X(f)\cap D_X(g)
\iff f(x)\neq0\text{ and }g(x)\neq0
\iff f(x)g(x)\neq0,
\]
so \(D_X(f)\cap D_X(g)=D_X(fg)\).

\paragraph{What happened mathematically.}
An open neighborhood of a point says that some closed system of equations fails there.
One failed equation already supplies a principal neighborhood.

\paragraph{Extensions.}
\begin{enumerate}
\item Show finite intersections of principal opens are principal.
\item Show every open cover can be refined by a cover of principal opens.
\item Revisit this argument in VI/08 for \(D(f)\subseteq\Spec A\).
\end{enumerate}

\paragraph{Provenance.}
Recovered from legacy DG4 affine Problem 3 and rewritten in canonical relative notation.
\fi
